% Generated draft: requires mathematical and source review.
\documentclass[11pt]{article}
\usepackage[margin=25mm]{geometry}
\usepackage{fontspec}
\setmainfont{Latin Modern Roman}
\usepackage{amsmath,amssymb,mathtools,unicode-math}
\setmathfont{Latin Modern Math}
\usepackage{enumitem,needspace,xurl,hyperref}
\setlength{\parindent}{0pt}
\setlength{\parskip}{5pt}
\setlength{\emergencystretch}{3em}
\newcommand{\N}{\mathbb{N}}
\newcommand{\Z}{\mathbb{Z}}
\newcommand{\Q}{\mathbb{Q}}
\newcommand{\R}{\mathbb{R}}
\newcommand{\C}{\mathbb{C}}
\newcommand{\F}{\mathbb{F}}
\newcommand{\A}{\mathbb{A}}
\newcommand{\PP}{\mathbb{P}}
\newcommand{\E}{\mathbb{E}}
\newcommand{\eps}{\varepsilon}
\newcommand{\dd}{\,\mathrm{d}}
\begin{document}
\section*{UnsolvedMath ID: 700664}
\textbf{Source problem code:} catalogue entry
\subsection{Definitions and mathematical statement}
Let \(G\) be a group with identity \(1\). A group is called simple if it is nontrivial and has no normal subgroups other than \(\{1\}\) and itself. For \(x\in G\), its conjugacy class is \(x^G=\{gxg^{-1}:g\in G\}\). For a subset \(X\subseteq G\) and a positive integer \(r\), define
\[
X^r=\{x_1\cdots x_r:x_1,\ldots,x_r\in X\}.
\]
Thus this notation means a product of exactly \(r\) factors, with repetitions allowed. An abelian subgroup \(H\leq G\) is maximal abelian if no strictly larger abelian subgroup of \(G\) contains it. Write
\[
N_G(H)=\{g\in G:gHg^{-1}=H\},\qquad C_G(h)=\{g\in G:gh=hg\}.
\]
The source asks whether there exists a group \(G\) satisfying the following conditions, with integers \(n,m\) interpreted as positive finite product lengths:
\begin{enumerate}
\item \(G\) is simple, and there exists \(n\geq 1\) such that
\[
G=(x^G)^n\qquad\text{for every }x\in G.
\]
The integer \(n\) is uniform over all conjugacy classes.
\item For any maximal abelian subgroups \(H,K\leq G\), there exists \(g\in G\) such that \(K=gHg^{-1}\).
\item For every maximal abelian subgroup \(H\leq G\),
\[
N_G(H)=H,\qquad C_G(h)=H\quad\text{for every }h\in H\setminus\{1\}.
\]
\item There exists \(m\geq 1\) such that, for every maximal abelian subgroup \(H\leq G\) and every \(a\in G\setminus H\),
\[
G=(aH)^m,
\qquad aH=\{ah:h\in H\}.
\]
Equivalently, for every such \(H,a\) and every \(g\in G\), there exist \(h_1,\ldots,h_m\in H\) with \(g=(ah_1)\cdots(ah_m)\). The integer \(m\) is uniform over \(H\), \(a\), and \(g\).
\end{enumerate}
No finiteness or infinitude assumption on \(G\) is supplied. There is an unresolved qualification in condition (1): the source says ``any conjugacy class,'' without excluding the identity class \(\{1\}\). The displayed formulation retains that literal quantifier. Its compatibility with the usual nontrivial meaning of simplicity requires clarification; an intended restriction to nontrivial conjugacy classes has not been silently inserted. The source also does not explicitly specify the signs of the integers \(n,m\); positive lengths are the standard interpretation of the stated products.
\subsection{Short English statement}
B. Poizat asks for a simple group with a uniform bound on products from each conjugacy class, conjugate maximal abelian subgroups that are self-normalizing and equal the centralizers of their nonidentity elements, and a uniform exact product length for every coset \(aH\) with \(a\notin H\). The source leaves the identity conjugacy class unexcluded, so the intended first condition needs clarification. A supplied 2005 comment reports a partial solution.
\subsection{Sources}
\begin{itemize}
\item Local UnsolvedMath registry, supplied complete source record: permanent ID \(700664\), exact source problem code \texttt{catalogue entry}, title \emph{Kourovka Notebook Problem 10.49}. The supplied registry status is \texttt{open}; this is a source label, not an independently verified mathematical or literature assessment.
\item \emph{The Kourovka Notebook (Unsolved Problems in Group Theory)}, original problem 10.49, 10th issue (1986); proposed by B. Poizat. The supplied background cites the 21st issue edition (2026), \texttt{arXiv:1401.0300v48}. The supplied collection description names E. I. Khukhro and V. D. Mazurov as editors.
\item Supplied \emph{Comment of 2005}: a partial solution is reported in E. Jaligot and A. Ould Houcine, \emph{J. Algebra} \textbf{280} (2004), 772--796. The supplied record gives neither an article title nor the scope of the partial solution. This historical report has not been independently verified here.
\end{itemize}
\end{document}
